\subsection{A genus-two quotient carrying the cubic comparison sign}

Use the finite-comparison notation
$U=\left(\begin{smallmatrix}1&1\\0&1\end{smallmatrix}\right)$,
$L=\left(\begin{smallmatrix}1&0\\-1&1\end{smallmatrix}\right)$,
$R=-UL$, $X=U^2$, $Y=L^2$, $C=5I$, and $Z=-I$ modulo eight.
Write $G_8=\mathrm{SL}_2(\mathbb Z/8)$ and
$P=\langle X,R\rangle$, of order $192$.
The already proved presentation has
\[
 N=\langle X,Y,C,Z\rangle,\quad |N|=64,\quad
 P=N\rtimes\langle R\rangle,
\]
with the unique forms $X^aY^bC^eZ^d$,
$0\leq a,b<4$, $0\leq e,d<2$, and
\[
 [X,Y]=C,\qquad RXR^{-1}=Y,\qquad
 RYR^{-1}=Z(XY)^{-1}.
\]
The elements $C,Z,X^2,Y^2$ are central in $N$.

In the target put
\[
 S=\begin{pmatrix}0&1\\-1&0\end{pmatrix},\quad
 T=\begin{pmatrix}1&1\\0&1\end{pmatrix},\quad
 A=I-S=\begin{pmatrix}1&2\\1&1\end{pmatrix}
 \quad\text{over }\mathbb F_3.
\]

\begin{proposition}\label{prop:genus-two-finite-quotient}
There is a surjection
\[
 \ell:G_8\longrightarrow\mathrm{GL}_2(\mathbb F_3),\qquad
 U\longmapsto A,\quad R\longmapsto T,
\]
whose kernel is the order-eight subgroup
\[
 K=\langle Z,X^2C,Y^2C\rangle.
\]
It kills $Z$, so the signed congruence representation induces
\[
 \bar\ell:\Gamma=C_2*C_3\longrightarrow
                 \mathrm{GL}_2(\mathbb F_3).
\]
For $s=uvu$, $r=uv$ and $u=r^{-1}s$, their images have orders
$2,3,8$, respectively.  On the subgroup $P$, this map agrees
with the previous interchange map on all commutator prowords in
$X,R$, although its value on $X$ is $S$ instead of $-S$.
\end{proposition}
\begin{proof}
The presentation of $P$ defines a character
$\chi:P\to\{\pm1\}$ by
\[
 \chi(X)=\chi(Y)=\chi(Z)=-1,\qquad
 \chi(C)=\chi(R)=1.
\]
For example the only nontrivial conjugation check is
$\chi(Y)=\chi(Z(XY)^{-1})=-1$; the remaining relations are
immediate.  Multiply the previously established interchange
homomorphism $\kappa:P\to\mathrm{SL}_2(\mathbb F_3)$ by this
central character.  The resulting surjection $\kappa'$ has
\[
 \kappa'(X)=S,\quad \kappa'(R)=T,\quad
 \kappa'(C)=-I,\quad\kappa'(Z)=I.
\]
Its kernel contains $K$.  The unique normal forms show that
$Z,X^2C,Y^2C$ are three independent commuting involutions, so
$|K|=8$.  Surjectivity and $|P|=192$ imply that this is the whole
kernel.  A character is trivial on commutators, giving the last
assertion about prowords.

The subgroup $P$ has index two in $G_8$, and $U^2=X\in P$.
The whole group is generated by $P$ and $U$: indeed
$L=U^{-1}ZR$.  Direct multiplication gives
\begin{equation}\label{eq:bolza-1}
UXU^{-1}=X,\qquad URU^{-1}=XYR^2.
\end{equation}
In the target, $A^2=S$, $A$ commutes with $S$, and
\begin{equation}\label{eq:bolza-2}
ATA^{-1}=S(TST^{-1})T^2.
\end{equation}
For clarity, both sides of \eqref{eq:bolza-2} are
$\left(\begin{smallmatrix}2&2\\1&0\end{smallmatrix}\right)$.
Since $X,R$ generate $P$, \eqref{eq:bolza-1}--\eqref{eq:bolza-2} show that conjugation by
$A$ on $\kappa'(P)$ agrees with conjugation by $U$ before
applying $\kappa'$.  Together with $A^2=\kappa'(U^2)$, these
identities give a well-defined homomorphism on the two cosets
$P\sqcup PU$, extending $\kappa'$ and sending $U$ to $A$.
Its image contains $\mathrm{SL}_2(\mathbb F_3)$ and the matrix
$A$ of determinant $-1$, hence is all of $\mathrm{GL}_2(\mathbb F_3)$.
Its kernel has order $384/48=8$ and is $K$.

The signed braid center maps to $Z$, so the map descends to
$\Gamma$.  There $r=uv$ maps to $T$, since $ZR$ maps to $T$;
and $s=ru$ maps to
\[
 TA=\begin{pmatrix}2&0\\1&1\end{pmatrix},
\]
an involution.  Finally $A^2=S$, $S^2=-I$ and $A^4=-I$,
so $A$ has order eight.  This proves the stated orders.
\end{proof}

\begin{corollary}\label{cor:genus-two-sign-cover}
The discrete kernel of $\bar\ell$ is the fundamental group of
a genus-two surface with six punctures.  Its compactification
carries a faithful deck action of $\mathrm{GL}_2(\mathbb F_3)$
of signature $(2,3,8)$.  The central involution is its
hyperelliptic involution.

On $D$ the induced action on this finite deck group is the
identity.  The pure harmonic conjugator maps to $I$, and the
pure cubic conjugator maps to $(-I)^{\tau(g)}$.
\end{corollary}
\begin{proof}
The subgroup has index $48$ in the modular group $\Gamma$.
It is torsion-free because the orders two and three of the
finite factors are preserved.  The cusp stabilizer has image
of order eight, so there are $48/8=6$ cusps.  Since
$\chi(\Gamma)=-1/6$, its Euler characteristic is $-8$.
Thus $2-2g-6=-8$ and $g=2$.  Filling the cusps gives the
signature $(2,3,8)$ action, by the same stabilizer counts.

The central matrix $-I=A^4$ lies in every conjugate of the
cusp stabilizer, so it fixes all six filled cusp points.
For an involution of a genus-two compact surface the
Riemann--Hurwitz formula is $2=4h-4+r$.  Since $r\geq6$,
one has $h=0$ and $r=6$.  This is the hyperelliptic involution.

On $D$ the established normalized signed covariance gives
$q_8\Phi_g=q_8$, hence identity on the deck group after $\ell$.
The established prime-two normalization gives
$q_8(\beta_g)=I$.  Also
$q_8(A_g)^{-1}=5^{\tau(g)}I=C^{\tau(g)}$, and
$\ell(C)=-I$.  Inversion leaves this sign unchanged, yielding
the asserted cubic conjugator value.
\end{proof}

This gives a concrete genus-two setting for the cubic finite
comparison. The equality $\tau=\eta$ requires the cubic branch-filling
comparison proved in Section~\ref{sec:normalized-cubic-identity}.

\paragraph{Identification with the classical Bolza cover.}
The genus-two cover in Proposition~\ref{prop:genus-two-finite-quotient}
is the classical Bolza cover.  Katz--Katz--Schein--Vishne
\cite[\S3 and Corollaries 13.3--13.4]{KKSV2016} describe its
index-two triangle-group inclusion and the corresponding
$\mathrm{SL}_2(\mathbb F_3)$ and $\mathrm{GL}_2(\mathbb F_3)$
quotients.  Their displayed matrices immediately before
Corollary 13.4 send the order-two and order-three generators to
\[
 s\longmapsto\begin{pmatrix}0&1\\1&0\end{pmatrix},\qquad
 r\longmapsto T=\begin{pmatrix}1&1\\0&1\end{pmatrix}.
\]
These are exactly conjugate to our pair:
\[
 T(TA)T^{-1}=\begin{pmatrix}0&1\\1&0\end{pmatrix},\qquad
 TTT^{-1}=T.
\]
Thus the kernels in the triangle group agree.  The compactification is the Bolza
curve $y^2=x^5-x$; the same model, its six hyperelliptic branch
points $0,\infty,\pm1,\pm i$, and its classical automorphism
group are recalled in \cite[\S4, equation (4.1)]{KRRBolza}.
That paper attributes the classification to Bolza's 1887 work.

Here we compute the factorization through
$\mathrm{SL}_2(\mathbb Z/8)$ with kernel
$\langle Z,X^2C,Y^2C\rangle$, together with the values of the
normalized frame conjugators. Section~\ref{sec:normalized-cubic-identity}
then identifies their cubic sign with $\eta$.
